    \documentclass{article}
    \usepackage{amsmath}
	\usepackage{amssymb}
    \usepackage[minted]{tcolorbox}
    \usepackage{xcolor} % Required for bgcolor in minted
    \definecolor{lightgray}{rgb}{0.9,0.9,0.9} % Define lightgray if not already defined
    \usepackage{fontspec} %fontspec is required for code font
    \setmonofont{JuliaMono}[Extension=.ttf, UprightFont=*-Regular, BoldFont=*-Bold, ItalicFont=*-RegularItalic, BoldItalicFont=*-BoldItalic, Contextuals=Alternate, Scale = 0.8]
    \usepackage{tabularx} % For tables that fit page width
    \usepackage{array} % For better column formatting
    \usepackage{booktabs} % For professional looking tables
	\usepackage{fullpage,graphicx,psfrag,amsfonts,verbatim, url}
    \usepackage{parskip}
    
    \title{\textbf{\textsf{\texttt{PEPit.jl} tutorial: Accelerated gradient method for smooth strongly convex minimization}}}
	
	\author{John Doe}
    
    \begin{document}
    \maketitle
    
    

\section{Setup}
Before we do anything let us load the necessary packages.

\begin{tcblisting}{listing only, minted language=julia, minted options={mathescape=true, breaklines}, colback=lightgray, colframe=black, boxrule=0pt, toprule=2pt, bottomrule=2pt, arc=0pt, outer arc=0pt, left=5pt, right=5pt, top=5pt, bottom=5pt}
using PEPit, OrderedCollections, Mosek, MosekTools, Clarabel
\end{tcblisting}


Let us start with defining an empty PEP, which we will construct step by step.

\begin{tcblisting}{listing only, minted language=julia, minted options={mathescape=true, breaklines}, colback=lightgray, colframe=black, boxrule=0pt, toprule=2pt, bottomrule=2pt, arc=0pt, outer arc=0pt, left=5pt, right=5pt, top=5pt, bottom=5pt}
problem = PEP()
\end{tcblisting}


\section{Problem in consideration}

Consider the convex minimization problem
\begin{equation*}
f_\star \triangleq \min_x f(x),
\end{equation*}
where $f$ is $L$-smooth and $\mu$-strongly convex. For this example, let us take $\mu=0.1$ and $L=1$. Let us declare this function type in \texttt{PEPit}. 

\begin{tcblisting}{listing only, minted language=julia, minted options={mathescape=true, breaklines}, colback=lightgray, colframe=black, boxrule=0pt, toprule=2pt, bottomrule=2pt, arc=0pt, outer arc=0pt, left=5pt, right=5pt, top=5pt, bottom=5pt}
μ = 0.1
L = 1
param = OrderedDict("mu" => μ, "L" => L)
\end{tcblisting}


This defines the set of parameters $\mu$ and $L$ to describe the function class. Okay, now time to define the function class itself.  

\begin{tcblisting}{listing only, minted language=julia, minted options={mathescape=true, breaklines}, colback=lightgray, colframe=black, boxrule=0pt, toprule=2pt, bottomrule=2pt, arc=0pt, outer arc=0pt, left=5pt, right=5pt, top=5pt, bottom=5pt}
func = declare_function!(problem, SmoothStronglyConvexFunction, param; reuse_gradient=true) # creates the class of strongly convex and smooth functions with parameters $\mu = 0.1$ and $L=1$
\end{tcblisting}


\section{Algorithm}

What we want to study in this tutorial is to compute the a worst-case guarantee for the  \textbf{fast gradient} method described as follows. 

\subsection{Algorithm:}
For $t \in \{0, \dots, n-1\}$,


\begin{equation*}
\begin{aligned}
    y_t & = x_t + \frac{\sqrt{L} - \sqrt{\mu}}{\sqrt{L} + \sqrt{\mu}}(x_t - x_{t-1}) \\
    x_{t+1} & = y_t - \frac{1}{L} \nabla f(y_t)
\end{aligned}
\end{equation*}
with $x_{-1}:= x_0$.

\section{Goal}

We want to compute the smallest possible $\tau(n, L, \mu)$ such that the guarantee
\begin{equation*}
f(x_n) - f_\star \leqslant \tau(n, L, \mu) \left(f(x_0) -  f(x_\star) + \frac{\mu}{2}\|x_0 - x_\star\|^2\right),
\end{equation*}
is valid, where $x_n$ is the output of the \textbf{fast gradient} method, and where $x_\star$ is the minimizer of $f$. Sometimes we say that $\left(f(x_0) -  f(x_\star) + \frac{\mu}{2}\|x_0 - x_\star\|^2\right)$ is the initial condition for this PEP. In short, for given values of $n$, $L$ and $\mu$,
$\tau(n, L, \mu)$ is computed as the worst-case value of
$f(x_n)-f_\star$ when $f(x_0) -  f(x_\star) + \frac{\mu}{2}\|x_0 - x_\star\|^2 \leqslant 1$.

\section{Construct the PEP step by step}

Let us implement the algorithm in PEPit now. For this intorductory example we will let $n=2$.

\begin{tcblisting}{listing only, minted language=julia, minted options={mathescape=true, breaklines}, colback=lightgray, colframe=black, boxrule=0pt, toprule=2pt, bottomrule=2pt, arc=0pt, outer arc=0pt, left=5pt, right=5pt, top=5pt, bottom=5pt}
n = 2
\end{tcblisting}


\subsection{Define $x_0, x_\star, f_\star$}
We start with defining the initial point $x_0 \equiv $ \texttt{x0}, opitmal point $x_\star \equiv$ \texttt{xs}, and the optimal objective value $f_\star \equiv$ \texttt{fs}.

\begin{tcblisting}{listing only, minted language=julia, minted options={mathescape=true, breaklines}, colback=lightgray, colframe=black, boxrule=0pt, toprule=2pt, bottomrule=2pt, arc=0pt, outer arc=0pt, left=5pt, right=5pt, top=5pt, bottom=5pt}
xs = stationary_point!(func) # creates $x_\star$

fs = value!(func, xs) # computes $f_\star = f(x_\star)$

x0 = set_initial_point!(problem) # creates $x_0$
\end{tcblisting}


\subsection{Define initial condition}

The next step is now define our initial condition, which is $f(x_0) -  f(x_\star) + \frac{\mu}{2}\|x_0 - x_\star\|^2 \leq 1$. 

\begin{tcblisting}{listing only, minted language=julia, minted options={mathescape=true, breaklines}, colback=lightgray, colframe=black, boxrule=0pt, toprule=2pt, bottomrule=2pt, arc=0pt, outer arc=0pt, left=5pt, right=5pt, top=5pt, bottom=5pt}
set_initial_condition!(problem, value!(func, x0) - fs + μ / 2 * (x0 - xs)^2 <= 1) # creates the initial condition $f(x_0) -  f(x_\star) + \frac{\mu}{2}\|x_0 - x_\star\|^2 \leq 1$
\end{tcblisting}


\subsection{Describe the algorithm}

Great, now it is time to implement the \emph{fast gradient method} in a loop. 

\begin{tcblisting}{listing only, minted language=julia, minted options={mathescape=true, breaklines}, colback=lightgray, colframe=black, boxrule=0pt, toprule=2pt, bottomrule=2pt, arc=0pt, outer arc=0pt, left=5pt, right=5pt, top=5pt, bottom=5pt}
κ = μ/L

x_new, y = x0, x0
for i in 1:n
    global x_new, y  # mark both as globals in this loop
    x_old = x_new
    x_new = y - 1 / L * gradient!(func, y)
    y = x_new + (1 - sqrt(κ)) / (1 + sqrt(κ)) * (x_new - x_old)
end
\end{tcblisting}


So when the loop ends \texttt{x\_new} contains $x_n$. Let us now evaluate $f(x_n)$.

\begin{tcblisting}{listing only, minted language=julia, minted options={mathescape=true, breaklines}, colback=lightgray, colframe=black, boxrule=0pt, toprule=2pt, bottomrule=2pt, arc=0pt, outer arc=0pt, left=5pt, right=5pt, top=5pt, bottom=5pt}
f_n = value!(func, x_new) # computes $f(x_n)$
\end{tcblisting}


Let us now tell the PEP that our performance measure is $f(x_N) - f_\star$.

\begin{tcblisting}{listing only, minted language=julia, minted options={mathescape=true, breaklines}, colback=lightgray, colframe=black, boxrule=0pt, toprule=2pt, bottomrule=2pt, arc=0pt, outer arc=0pt, left=5pt, right=5pt, top=5pt, bottom=5pt}
set_performance_metric!(problem, value!(func, x_new) - fs) # set performance metric to $f(x_n) - f_\star$
\end{tcblisting}


\subsection{Solve the PEP}
We have everything to build up the PEP. Now time to solve!

\begin{tcblisting}{listing only, minted language=julia, minted options={mathescape=true, breaklines}, colback=lightgray, colframe=black, boxrule=0pt, toprule=2pt, bottomrule=2pt, arc=0pt, outer arc=0pt, left=5pt, right=5pt, top=5pt, bottom=5pt}
τ_PEPit = solve!(problem; solver = Mosek.Optimizer, verbose=true) # we can use other SDP solvers such as Clarabel, via `solver = Clarabel.Optimizer`
\end{tcblisting}


\subsection{Comparison with theoretical results}

If everything went well, we should get \texttt{τ\_PEPit ≈ 0.34760221803672986} or something very close to it. What does the theory say in this regard? The following \textbf{upper} guarantee can be found in [1, Corollary 4.15]:

\begin{equation*}
f(x_n)-f_\star \leqslant \left(1 - \sqrt{\frac{\mu}{L}}\right)^n \left(f(x_0) -  f(x_\star) + \frac{\mu}{2}\|x_0 - x_\star\|^2\right).
\end{equation*}

Let us compute the theoretical value of $\tau$.

\begin{tcblisting}{listing only, minted language=julia, minted options={mathescape=true, breaklines}, colback=lightgray, colframe=black, boxrule=0pt, toprule=2pt, bottomrule=2pt, arc=0pt, outer arc=0pt, left=5pt, right=5pt, top=5pt, bottom=5pt}
τ_Theory = (1 - sqrt(κ))^n

@info "💻  PEPit guarantee: f(x_n)-f_*  <= $(round(τ_PEPit, digits=6)) (f(x_0) - f(x_*) + mu/2*||x_0 - x_*||^2)"

@info "📝 Theoretical guarantee: f(x_n)-f_*  <= $(round(τ_Theory, digits=6)) (f(x_0) - f(x_*) + mu/2*||x_0 - x_*||^2)"
\end{tcblisting}


So the gurantee that we found via \texttt{PEPit} is tighter than the theoretical guarantee!

\textbf{References}:

[1] A. d’Aspremont, D. Scieur, A. Taylor (2021). Acceleration Methods. Foundations and Trends
in Optimization: Vol. 5, No. 1-2.


\end{document}
